% Faithful transcription — original Latin. Transcribed from C. I. Gerhardt (ed.), Leibnizens
% mathematische Schriften (Internet Archive scan leibnizensmathe07leibgoog), printed pages
% 234–237 (= scan leaves n250–n253): Leibniz's "De linea isochrona…" (Acta Eruditorum, 1689).
% \origpage uses the PRINTED page numbers. Only Leibniz's text is transcribed: Gerhardt's
% editorial footnotes (in German) and the appended "Beilage" (Huygens's solution, in French) are
% NOT part of this work and are omitted. Multi-letter point/curve labels stay plain math letters.
\documentclass{article}
\usepackage{readmasters}
\begin{document}

\origpage{234}
\section*{III. De linea isochrona, in qua grave sine acceleratione descendit, et de controversia cum Dn. Abbate De Conti}

Cum a me in his Actis Martio 1686 editis publicata esset demonstratio contra \emph{Cartesianos},
qua vera virium aestimatio traditur, ostenditurque non quantitatem motus, sed potentiae, a
quantitate motus differentem servari, Vir quidam doctus in Gallia, Dn. Abbas \emph{De Conti}, pro
Cartesianis respondit, sed, ut post apparuit, vi mei argumenti non satis perspecta. Credidit enim,
recepta quaedam alia principia a me impugnari, quae in Novellis Reip. Litterar. mens. Jun. 1687
p.~579 enumerat, et negat p.~519 seq. se agnoscere contradictionem, quam ego in illis invenire mihi
videar: cum tamen nunquam mihi de illis dubitare in mentem venerit, quemadmodum ipsum admonui
Novell. Reip. Lit. Septemb. 1687. Idem ut eluderet objectionem meam, conjecerat se in diverticulum
temporis, quod eo modo, quo conceptus a me erat status controversiae, plane est accidentale. Eadem
enim manente altitudine, eadem vis acquiritur aut impenditur a gravibus quocunque tempore indulto,
quod pro inclinatione descensus majore minoreve augetur aut minuitur. Ea occasione, quo magis
appareret, tempus atque adeo distinctionem inter potentias isochronas vel anisochronas hoc loco
nihil ad rem facere, et ut ex disputatione nostra aliquid incrementi scientia caperet,
\emph{problema} tale, a me inter scribendum solutum, et, ut videtur, non inelegans, ipsi proposui
in dictis Novellis Septembr. 1687: ``Invenire lineam isochronam, in qua grave descendat
uniformiter, sive aequalibus temporibus aequaliter accedat ad horizontem, atque adeo sine
acceleratione et aequali semper velocitate deorsum feratur.'' Sed Dn. Abbas \emph{De Conti} nihil
ultra reposuit, sive quod problema attingere nollet, sive quod agnita tandem mente mea,
satisfactum sibi judicaret. Sed ejus loco problema hoc sua opera dignum judicavit Vir celeberrimus
\emph{Christianus Hugenius}, cujus solutio mea prorsus consona extat in Novellis Reip. Lit. Octbr.
1687, sed suppressa demonstratione et explicatione
\origpage{235}
discriminis inter diversas lineas ejusdem, ut ait, generis, quas satisfacere notat. Haec igitur
ego supplere hoc loco volui, facturus citius, nisi aliquid hic a Dn. Abbatis industria
exspectavissem.

\emph{Problema.} Invenire lineam planam, in qua grave sine acceleratione descendit.

\emph{Solutio.} Sit (fig.~116)\ednote{The plate containing fig.~116 could not be located in the
scans available for this work (the Gerhardt edition scan has no plates section); it is not
reproduced here.} linea paraboloeides quadrato-cubica quaecunque $\beta Ne$ (nempe ubi
solidum sub quadrato basis $NM$ et parametro $aP$ aequale est cubo altitudinis $\beta M$) ita sita,
ut verticis $\beta$ tangens $\beta M$ sit perpendicularis horizonti, in cujus lineae puncto
quocunque $N$ si ponatur grave ea descendendi ulterius celeritate praeditum, quam potuit acquirere
descendendo ex horizonte $Aa$, cujus elevatio $a\beta$ supra verticem $\beta$ sit
$\frac{4}{9}$ parametri curvae, tunc idem grave descendet porro uniformiter per lineam $Ne$,
utcunque continuatam, ut desiderabatur.

\emph{Demonstratio.} Recta $NT$ curvam $\beta Ne$ tangat in $N$ et ipsi $\beta M$ occurrat in $T$.
Utique (ex nota proprietate tangentium hujus curvae) erit $TM$ ad $NM$ in subduplicata ratione
$a\beta$ ad $\beta M$. Ergo $TM$ ad $TN$ erit in subduplicata ratione $a\beta$ ad
$a\beta + \beta M$ seu ad $aM$. Jam ratio $TM$ ad $TN$ eadem est, quae velocitatis per curvam
descendendi (seu horizonti porro in curva appropinquandi), quam grave habet positum in $N$ ad
velocitatem, qua idem ex $N$ porro, non per curvam, sed libere descenderet, si posset (ut constat
ex natura motus inclinati). Sed velocitas haec libera porro est ad constantem quandam in
subduplicata ratione $aM$ ad $a\beta$; sunt enim (ut ex motu gravium constat) velocitates liberae
in altitudinum
\origpage{236}
(unde descendendo quaesitae sunt) $aM$ subduplicata ratione: ergo velocitas descendendi per curvam
$Ne$, quam grave habet in quocunque curvae puncto $N$ positum, est ad velocitatem constantem in
composita subduplicata ratione $a\beta$ ad $aM$ et $aM$ ad $a\beta$, quae est ratio aequalitatis.
Ipsamet igitur velocitas illa per curvam descendendi est constans, seu ubique in curva $Ne$ eadem.
Quod praestandum erat.

\emph{Consectaria:} 1) Grave celeritatem habens tamquam lapsum ab altitudine aliqua $Aa$,
descendere potest ex eodem puncto $N$ per curvas isochronas infinitas, sed ejusdem speciei, seu
sola magnitudine parametri differentes, ut $Ne$, $N(e)$, $NE$, quae omnes sunt paraboloeides
quadrato-cubicae, adeoque similes inter se. Imo quaelibet hujusmodi paraboloeidum hic inservit,
modo ita collocetur, ut $a\beta$ vel $(a)(\beta)$ distantia verticis ab horizontali $a(a)$, unde
descendere incepit grave, sit $\frac{4}{9}$ parametri curvae $\beta e$ vel $(\beta)(e)$: nec refert,
an grave isochrone descensurum in curva $N(e)$ pervenerit ad $N$ ex $a(a)$ per viam
$(a)(\beta)N$, an per aliquam aliam, aut sine descensu ullo ob aliam causam eandem celeritatem
atque directionem acquisiverit. Ex infinitis tamen istis lineis isochronis, in quibus grave ex $N$
porro sine acceleratione descendere potest, ea celerrimum ipsi descensum praebet, cujus vertex est
ipsum punctum $N$, qualis est $NE$, quam recta $AN$ horizonti perpendicularis tangit.

2) Tempus descensus per rectam $a\beta$ est ad tempus descensus per curvam $\beta N$, ut dimidia
altitudo $\beta M$ ad ipsam $a\beta$: ac proinde si $\beta M$ sit dupla $a\beta$, aequalia erunt
tempora descensuum per $a\beta$ et per $\beta N$. Quorum ratio manifesta est: nam tempora descensus
uniformis sunt inter se ut altitudines, et ex demonstratis a \emph{Galilaeo}, tempus quo mobile
percurrit altitudinem $a\beta$ motu accelerato, est duplum ejus, quo percurrit aequalem altitudinem
$\beta M$ (ut hoc loco fit, licet per curvam $\beta N$) motu uniformi, qui celeritatem habet
aequalem ultimae per accelerationem acquisitae in $\beta$.

Hoc autem problema fateor me non Geometris primariis proposuisse, qui interiorem quandam Analysin
callent, sed his potius, qui cum \emph{erudito illo Gallo} sentiunt, quem mea de \emph{Cartesianis}
plerisque hodiernis (Magistri paraphrastis potius, quam aemulatoribus) querela suboffendisse
videbatur. Tales enim cum alias receptis inter Cartesianos dogmatibus, tum etiam analysi inter
ipsos pervulgatae nimium tribuunt, adeo ut se ipsius ope quidvis in Mathesi
\origpage{237}
(si modo velint scilicet calculandi laborem sumere) praestare posse arbitrentur, non sine
detrimento scientiarum, quae falsa jam inventorum fiducia negligentius excoluntur. His materiam
exercendae suae Analyseos praebere volueram in hoc problemate, quod non prolixo calculo, sed arte
indiget.

Si quis tamen praereptam sibi jam solutionem queratur, poterit \emph{aliam isochronam} huic
vicinam quaerere, in qua non, ut hactenus, grave uniformiter recedat ab horizontali (vel ad eam
accedat), sed a certo puncto. Unde problema erit tale, \emph{invenire lineam, in qua descendens
grave recedat uniformiter a puncto dato, vel ad ipsam accedat.}

Talis foret linea $NQR$, si ejus esset naturae, ut ex puncto dato seu fixo $A$, ductis rectis
quibuscunque ad curvam ut $AN$, $AQ$, $AR$, esset excessus $AR$ super $AQ$, ad excessum $AQ$ super
$AN$, ut tempus quo descenditur per arcum $QR$, ad tempus quo descenditur per arcum $NQ$.

\end{document}
